\section{Operational Resilience and Learning}

The first thirteen sections are largely about preventing or catching an error before or as it happens. This section covers what a complete framework also needs: what happens to state that persists after an output is accepted, how the human and the system actually perform as a joint unit rather than as a producer plus a separate checker, how success at one level of evaluation does or doesn't transfer to the next, the ordinary (non-adversarial) software failures that can look like model failures but aren't, and what happens after an error has already escaped.

\textbf{Workflow state as its own risk surface, distinct from the model or the specific output.} An incorrect output doesn't necessarily stay contained to the interaction that produced it --- it can become a matter record, a precedent entry, a knowledge-management artifact, a future retrieval result, a worked example folded into someone else's prompt, evaluation data, fine-tuning data, or an input to an entirely separate automated decision downstream. Once any of that happens, the system isn't merely at risk of reproducing the same error again --- it is, in a real sense, \textbf{learning organizationally from its own unverified output}, the way the worked example in Section 12 ends up doing. This is likely to matter more over time than isolated single-output hallucinations, precisely because persistent memory and organizational knowledge bases are becoming a standard part of how these systems are deployed. Worth building in as controls: keeping generated material visibly distinguished from authoritative material; provenance that survives copying and summarization rather than getting stripped along the way; a quarantine period before AI-produced content is allowed into a trusted corpus; a promotion gate requiring actual approval before that happens; expiration and revalidation dates on anything that was promoted; a rule against letting multiple generated claims be treated as independent corroboration of each other, since they may share the same underlying source or the same underlying error; the ability to trace any derivative artifact back to its original evidence; and a process for correcting or removing contaminated downstream records once a source error is found. This is, in effect, ordinary data-integrity management applied to AI-generated content specifically.

\textbf{Modeling the human-AI team as a joint system, not as "the model, plus a reviewer" bolted on afterward.} Section 9 already names complacency as a failure mode; the fuller version treats the human reviewer as another fallible subsystem whose behavior actively changes in response to the automation around it, rather than as a constant, independent check. The relevant unit to evaluate is not "model accuracy" plus "human review" added together as if they were separate and additive --- it's \textbf{joint human-system performance under realistic workload, interface, and incentive conditions}, which can be considerably worse (or occasionally better) than either component considered alone. Worth asking directly: does the model's answer anchor the reviewer's own initial view before they've evaluated it independently? Is disagreeing with the system's output psychologically or procedurally costly in a way that discourages it? Does the interface actually make contrary evidence easy to inspect, or does it bury it? Can a reviewer even tell which parts of an output were actually checked versus merely displayed? Does performance change noticeably after a long run of outputs that all turned out to be fine --- this is the same vigilance-decay mechanism named in Section 9, and asking specifically how performance shifts after a run of consecutive acceptable outputs is a useful \emph{kind} of question to ask in a given deployment, though any specific number used to illustrate it should be read as illustrative rather than as an established general threshold. Does higher throughput reduce the time actually available per review? Does reliance on the system erode the reviewer's own independent skill over time? And, less obviously but importantly: who does each party --- the reviewer and the organization --- actually believe is responsible if something goes wrong? Worth measuring where possible: human-alone performance, AI-alone performance, and combined performance as three separate numbers rather than assuming the third is simply derived from the first two; the reviewer's override rate, split into correct and incorrect overrides; the detection rate for errors deliberately seeded into a review batch as a check on the check; and review time and fatigue effects across a session.

\textbf{Separating evaluation levels, since success at one doesn't establish success at the next.} A useful hierarchy: \textbf{model evaluation} (what can the underlying model do, in isolation); \textbf{component evaluation} (do retrieval, parsing, and individual tools work correctly on their own); \textbf{system evaluation} (does the assembled product work end to end); \textbf{workflow evaluation} (does it work with actual people, actual policies, and actual workloads, not a clean test scenario); and \textbf{field evaluation} (what actually happens after sustained real-world use, once novelty and careful onboarding attention have worn off). A model can pass a benchmark while the product built around it retrieves the wrong documents in practice. A product can pass every scripted test while real users misunderstand the interface in ways no test scenario anticipated. A pilot can work cleanly while production volume later creates time pressure, shortcuts, and queue-driven behavior that a small pilot never surfaced. This staged distinction has a real institutional analogue: NIST's ARIA program, a three-tier evaluation approach launched to assess large language model applications, explicitly separates model testing, red-teaming, and field testing specifically because it found that traditional accuracy benchmarks don't capture how a system behaves once real people are using it in realistic settings --- model testing is necessary but was explicitly found insufficient on its own. A thorough evaluation, at whichever level, should include representative task sampling (not just favorable or convenient examples), rare and high-consequence cases specifically, subgroup and scenario breakdowns rather than only an aggregate score, confidence intervals rather than a single point estimate, a measure of adjudicator disagreement where human judgment is involved, an explicit list of what was excluded from the evaluation and why, challenge cases not designed by the product team itself, and field observation genuinely after deployment, not only before it.

\textbf{Conventional software and distributed-system failures, since the AI frame can make every failure look like a model failure when it isn't one.} An AI-enabled workflow can fail for reasons that have nothing to do with the model at all: stale data, API timeouts, incorrect retry logic, partial transactions, duplicated actions, race conditions, permission misconfiguration, a broken parser, silent truncation of a long input or output, malformed tool output, incorrect version routing between model updates, failed logging that leaves no trace of what actually happened, or an unavailable rollback mechanism when something needs undoing. A model can make a perfectly correct decision and still produce the wrong effect if the resulting action executes twice, gets applied against stale state, or only partially completes. The action-and-effect function named in Section 4, and the containment reasoning in Section 6, both benefit from a complementary, non-adversarial version of the same caution: is the action idempotent (does running it twice do the same thing as running it once)? Is it transactional (does it either complete fully or not at all, with no partial state left behind)? Is completion actually confirmed independently, rather than merely assumed? Can partial success be detected as distinct from full success or full failure? Is the state the model actually observed and reasoned about still the state that's true at the moment the action executes, or has it gone stale in between? And can retrying after an apparent failure make the outcome worse rather than better? This reinforces a point made throughout this report: the object that actually needs evaluating is the whole workflow, not the model considered in isolation.

\textbf{Resilience and learning after an error has already escaped, since pre-action verification is only one layer of defense, not the whole of it.} Worth having in place: graceful degradation rather than an all-or-nothing failure; circuit breakers that halt a process automatically past a defined threshold; transaction limits that cap how much damage a single run can do; a working rollback path; an out-of-band shutdown mechanism that doesn't depend on the system's own cooperation to function --- echoing the containment reasoning in Section 6 directly; canary and shadow deployments that surface problems before they reach full production; predefined stop conditions agreed in advance, rather than decided under pressure in the moment; a real incident-response process; a way to capture and learn from near misses, not only from actual incidents; structured post-incident causal analysis; periodic revalidation on a schedule rather than only when something visibly breaks; and named triggers that automatically invalidate a prior approval --- a change to the model, the prompt, the retrieval corpus, the permission structure, the available tool set, the user population, the interface, the workload, or the applicable law should all be treated as events that expire an existing sign-off, not as background changes an old approval silently continues to cover. This connects directly back to Section 3's requirement that every assessment terminate in a conditional decision: that decision should carry an explicit expiry condition, not just an approval date.

\textbf{A practical worksheet, useful for making this operational rather than only conceptual --- with one caveat about how to use it.} For a consequential node or handoff in a workflow, it's worth recording: the intended transformation (what this step should change or decide); input assumptions (what has to be true for it to work); failure modes (what can be wrong, omitted, stale, or unauthorized here); the hazard (what unsafe state could result); propagation (where that state travels next); load-bearing dependencies (which later conclusions actually rely on this one); the verification mechanism and what specific failure it targets; its independence from whatever it's checking; its coverage, and just as importantly what it doesn't examine; the consequence profile (severity, blast radius, persistence); detectability (how and when the failure would actually be discovered); recovery options (can it be stopped, corrected, or reversed); the ongoing monitoring signal that shows continued performance; the named owner who accepts residual risk; the stop condition that would suspend or roll back the workflow; and an honest list of what simply hasn't been established yet. This worksheet should be applied with the same proportionality Section 3 already establishes for the overall framework --- filling in all sixteen fields with genuine rigor for every node of a low-stakes internal brainstorming tool would be disproportionate effort in the wrong direction, and doing so risks producing exactly the false precision the rest of this section is trying to guard against: a sensitivity or specificity figure that looks measured but was actually estimated on the spot is worse than an honest "not established," because it invites the same over-trust in a checker that Section 9 already warns against for model confidence. Reserve full rigor for the nodes where the hazard analysis in Section 11 actually indicates the stakes justify it.

In short, this section and the three before it add four questions to the framework's earlier ones: what outcome is actually being produced, and what losses specifically must be prevented; what is known, uncertain, disputed, or missing at each point; where a faulty state can enter the workflow and how it propagates or becomes embedded once it does; and what evidence could detect it, how independent that evidence actually is, what would contain or reverse the consequences, and how the organization will know when its underlying assumptions have stopped holding true.
